% !TeX root = ../main.tex

% ==================================================
% 1.8. Process and Program
% ==================================================



\subsection{프로세스와 외부 프로그램}

\subsubsection{프로그램(Program)과 프로세스(Process)}

가장 먼저 구분해야 할 개념은 프로그램과 프로세스이다.

다음 파일을 생각해 보자.

\begin{dhkcode}
# train.py
print("Training started")
\end{dhkcode}

이 파일은 디스크에 저장된 Python 소스 코드이다.

아직 실행하지 않았으므로 \zsh{train.py} 자체는 프로세스가 아니다.

이제 터미널에서 실행해 보자.

\begin{dhkzsh}
python train.py
\end{dhkzsh}

이 명령이 실행되는 과정을 단순화하면 다음과 같다.

\begin{dhkbranch}[in=8]
node main :: Terminal
arrow main
node main :: Shell
arrow main :: python train.py
node main :: Operating System
arrow main :: Python 실행 요청
node main :: Python Process
arrow main :: Python Interpreter
node main :: train.py 실행
\end{dhkbranch}

여기서 중요한 점은 다음과 같다.

\begin{dhksmry}
    - 프로그램(Program) : 실행 가능한 명령이나 코드가 저장된 것

    - 프로세스(Process) : 운영체제에서 실행 중인 프로그램의 인스턴스
\end{dhksmry}

따라서 \zsh{train.py}는 소스 파일이고,

이를 실행하는 Python 인터프리터가 프로세스로 동작하는 것이다.

정확하게 말하면

\zsh{python train.py}는 Python 인터프리터 프로그램을 실행하고,

그 인터프리터가 \zsh{train.py}를 읽어서 실행하는 구조이다.

같은 프로그램을 두 번 실행하면?

터미널을 두 개 열어서 각각 실행한다고 생각해 보자.

\begin{dhkzsh}
python train.py
\end{dhkzsh}

\begin{dhkzsh}
python train.py
\end{dhkzsh}

그러면 일반적으로 두 개의 독립적인 Python 프로세스가 생성된다.

\begin{dhkdir}
Python Interpreter
> Process A
>> train.py
> Process B
>> train.py
\end{dhkdir}

두 프로세스는 같은 소스 파일을 실행하지만,

서로 다른 실행 상태를 가지고 있다.

예를 들어 각각 독립적인 변수와 Python 객체를 생성할 수 있다.

\begin{dhkcode}
# train.py

epochs = 10
epochs += 1

print(epochs)
\end{dhkcode}

두 프로세스에서 실행해도 각각 11이 출력된다.

왜냐하면 두 프로세스의 \zsh{epochs}가

동일한 Python 객체를 공유하는 것이 아니기 때문이다.

물론 프로세스끼리 데이터를 주고받는 방법도 있다.

하지만 이를 위해서는 별도의 통신 또는 공유 메커니즘이 필요하다.

\subsubsection{프로세스에는 무엇이 존재할까?}

프로세스를 단순히 실행 중인 코드라고 생각하면 중요한 부분을 놓칠 수 있다.

운영체제는 프로세스를 관리하기 위해 여러 정보를 유지한다.

대표적으로 다음과 같은 것들이 있다.

\begin{dhktable}
|구성 요소|설명|
|:---|:---|
|PID|프로세스를 식별하는 번호|
|Memory|프로세스가 사용하는 메모리 공간|
|Environment Variables|프로세스가 전달받은 환경변수|
|CWD|현재 작업 디렉터리|
|Arguments|실행할 때 전달받은 인수|
|stdin|표준 입력|
|stdout|표준 출력|
|stderr|표준 오류 출력|
|Exit Code|프로세스의 종료 상태|
\end{dhktable}

익숙한 개념들이 보인다.

우리는 이미 다음을 배웠다.

\begin{dhkcode}
import os
import sys
from pathlib import Path

print(os.environ)
print(Path.cwd())
print(sys.argv)
print(sys.executable)
\end{dhkcode}

이 코드들이 서로 무관한 기능을 제공하는 것이 아니다.

현재 실행 중인
\begin{dhksmry}[centered]
\textbf{Python 프로세스의 실행 환경과 상태를 서로 다른 관점에서 조회하는 것}
\end{dhksmry}

이다.

\zsh{sys.executable}은 현재 Python 인터프리터의 실행 파일 경로를 나타내고,

\zsh{sys.argv}는 Python 프로그램에 전달된 인수를 나타낸다는 차이가 있다.

기존에 배웠던 개념들이 하나의 프로세스라는 단위로 연결되기 시작한다.

\subsubsection{\zsh{os}와 \zsh{sys}의 차이}

지금까지 \zsh{os}와 \zsh{sys}를 여러 번 사용했다.

그런데 둘의 역할이 비슷해 보여서 오해할 수 있다.

\begin{dhknote}
\begin{dhksmry}[centered]
sys - Python 실행 환경과 인터프리터
\end{dhksmry}

\begin{dhksmry}
sys는 Python 인터프리터와 실행 환경에 관련된 기능을 제공한다.
\end{dhksmry}
\end{dhknote}

\begin{dhkcode}
import sys

print(sys.executable)
print(sys.version)
print(sys.platform)
print(sys.argv)
print(sys.path)
\end{dhkcode}

각각 다음을 의미한다.

\begin{dhktable}
|속성|의미|
|:---|:---|
|sys.executable|현재 인터프리터 실행 파일 경로|
|sys.version|Python 버전 정보|
|sys.platform|실행 플랫폼 식별|
|sys.argv|커맨드라인 인수|
|sys.path|모듈 탐색 경로|
\end{dhktable}

\begin{dhknote}
\begin{dhksmry}[centered]
os - 운영체제와 상호작용
\end{dhksmry}

\begin{dhksmry}
os는 운영체제가 제공하는 기능과 상호작용하기 위한 인터페이스다.
\end{dhksmry}
\end{dhknote}

\begin{dhkcode}
import os

print(os.getcwd())
print(os.environ)
print(os.getpid())
\end{dhkcode}

\begin{dhktable}
|함수, 속성|의미|
|:---|:---|
|os.getcwd()|현재 작업 디렉터리|
|os.environ|현재 프로세스의 환경변수|
|os.getpid()|현재 프로세스의 PID|
\end{dhktable}

예를 들어 다음 코드를 실행해 보자.

\begin{dhkcode}
import os

print(os.getpid())
\end{dhkcode}

출력은:

\begin{dhkzsh}
42817
\end{dhkzsh}

와 같을 수 있다.

여기서 \zsh{42817}은 운영체제가 해당 프로세스에 부여한 PID이다.

같은 코드를 새 Python 프로세스에서 실행하면 일반적으로 다른 PID가 출력된다.

그렇다면 왜 \zsh{os}와 \zsh{sys}가 분리되어 있을까?

역할을 구분하기 위해서이다.

\begin{dhkdir}
Python Program
> sys
>> Python Interpreter 관련 기능
> os
>> Operating System 관련 기능
\end{dhkdir}

물론 두 모듈의 기능이 완전히 배타적인 것은 아니다.

예를 들어 둘 다 프로세스 종료와 관련된 기능을 제공한다.

따라서 이렇게 이해하는 것이 좋다.

\begin{dhksmry}[centered]
sys는 \textbf{Python 인터프리터 관점}

os는 \textbf{운영체제 인터페이스 관점}
\end{dhksmry}

에 초점을 맞춘 모듈이다.

\subsubsection{부모 프로세스와 자식 프로세스}

여기서부터는 이전에 배운 환경변수와 직접 연결된다.

터미널에서 다음을 실행했다고 생각해 보자.

\begin{dhkzsh}
python train.py
\end{dhkzsh}

일반적인 터미널 환경에서는 Shell이 Python 프로세스를 실행한다.

\begin{dhkbranch}
node main :: Shell Process
arrow main :: python train.py
node main :: Python Process
line main
node main :: train.py
\end{dhkbranch}

이때 \textbf{Shell을 부모 프로세스(Parent Process)},

Python을 \textbf{자식 프로세스(Child Process)}라고 부를 수 있다.

자식 프로세스는 생성될 때 부모가 설정해 전달한 환경변수 등을 받는다.

예를 들어 macOS에서 다음과 같이 실행하면:

\begin{dhkzsh}
export MODEL_NAME=bert-base-uncased

python train.py
\end{dhkzsh}

Python 코드에서 다음을 읽을 수 있다.

\begin{dhkcode}
import os

print(os.environ.get("MODEL_NAME"))
\end{dhkcode}

출력:

\begin{dhkzsh}
bert-base-uncased
\end{dhkzsh}

이전에 배웠던 환경변수 전달이 실제로는 이러한 프로세스 생성 과정에서 이루어진다.

단, 환경변수가 부모 프로세스와 자식 프로세스 사이에서 자동으로 실시간 동기화되는 것은 아니다.

자식 프로세스는 생성될 때 전달받은 환경을 바탕으로 자신의 환경변수를 관리한다.

\subsubsection{Python에서 다른 프로그램을 실행할 수 있을까?}

가능하다.

지금까지는 Shell에서 Python을 실행했다.

\begin{dhkdir}
Shell
> Python
\end{dhkdir}

그런데 Python에서도 새로운 프로세스를 실행할 수 있다.

\begin{dhkdir}
Shell
> Python
>> External Program
\end{dhkdir}

이때 대표적으로 사용하는 모듈이 바로 \zsh{subprocess}이다.

예를 들어:

\begin{dhkcode}
import subprocess

subprocess.run(["python", "--version"])
\end{dhkcode}

과 같은 코드가 있다.

이 코드는 Python 프로그램 안에서 외부 명령을 실행한다.

여기서 \zsh{subprocess.run()}에 전달한 리스트를 살펴보자.

\begin{dhkzsh}
["python", "--version"]
\end{dhkzsh}

\begin{dhksmry}
    - "python" : 실행할 프로그램

    - "--version" : 해당 프로그램에 전달할 인수
\end{dhksmry}

즉, 터미널에서 실행하는:

\begin{dhkzsh}
python --version
\end{dhkzsh}

과 비슷한 작업을 Python 코드 안에서 수행하는 것이다.

다만 중요한 차이가 있다.

\zsh{subprocess.run()}은 기본적으로 Shell을 거치지 않고 프로그램을 직접 실행한다.

그리고 \zsh{"python"}이라는 이름은 실행 환경의 프로그램 검색 규칙에 따라 해석되므로,

반드시 현재 Python 프로세스와 같은 인터프리터를 실행한다고 보장할 수 없다.

현재 사용 중인 Python 인터프리터를 정확하게 실행하고 싶다면 다음처럼 작성하는 것이 좋다.

\begin{dhkcode}
import subprocess
import sys

subprocess.run([sys.executable, "--version"])
\end{dhkcode}

이렇게 하면 \zsh{sys.executable}이 가리키는 인터프리터를 사용하게 된다.

특히 Conda 환경을 사용한다면 중요한 부분이다.

\subsubsection{\zsh{subprocess.run()}의 기본 실행 흐름}

다음 코드를 보자:

\begin{dhkcode}
import subprocess
import sys

print("A")

subprocess.run([sys.executable, "--version"])

print("B")
\end{dhkcode}

실행 흐름은 다음과 같다.

\begin{dhkbranch}[in=8]
node main ::Parent Python Process
branchline main -> pr1 :: print("A")
stop pr1
branchline main -> spr :: subprocess.run()
arrow spr
node spr :: Child Process //* python --version
stop spr
handoff main -> cpe :: Child Process 종료
arrow cpe
node cpe :: print("B")
\end{dhkbranch}

일반적인 \zsh{subprocess.run()}은 자식 프로세스가 종료될 때까지 기다린다.

따라서 출력 순서는 다음과 같다.

\begin{dhkzsh}
A
Python 3.x.x
B
\end{dhkzsh}

단, Jupyter 등에서 실행하면 출력 캡쳐 방식에 따라 화면에 표시되는 순서는 다르게 보일 수 있다.

여기서 가장 중요한 개념은 다음이다.

\begin{dhknote}
\begin{dhksmry}[centered]
\zsh{subprocess.run()}은 외부 프로그램을 실행하고,

해당 프로그램이 종료될 때까지 기다린 뒤 결과를 반환한다.
\end{dhksmry}
\end{dhknote}

반대로 자식 프로세스를 실행한 뒤 종료를 기다리지 않고

다른 작업을 진행하고 싶다면 \zsh{subprocess.Popen()} 등을 사용할 수 있다.

\textbf{여기까지의 핵심 연결}

\begin{dhkbranch}[in=8]
node main :: Operating System
arrow main
node main :: Shell Process
arrow main :: python train.py --epochs 10
node main :: Python Process
branchline main -> pi :: Python Interpreter //* sys.executable
stop pi
branchline main -> cla :: Command-line Arguments //* sys.argv
stop cla
branchline main -> ev :: Environment Variables //* os.environ
stop ev
branchline main -> cwd :: Current Working Directory //* Path.cwd()
stop cwd
branchline main -> pid :: Process ID //* os.getpid()
stop pid
handoff main -> spr :: subprocess.run()
arrow spr
node spr :: Child Process
\end{dhkbranch}

이 구조가 8장의 출발점이다.

\textbf{앞으로 배울 내용}

\begin{dhksmry}
    1. subprocess.run()의 반환값과 CompletedProcess
    
    2. stdin, stdout, stderr의 의미
    
    3. returncode와 exit code
    
    4. sys.exit()와 os._exit()의 차이
    
    5. Shell의 파이프 |와 리다이렉션 >
    
    6. subprocess.run()에서 출력 캡처하기
    
    7. shell=True의 의미와 주의점
    
    8. subprocess.Popen()과 비동기적 프로세스 실행
    
    9. ML/DL 프로젝트에서 외부 학습 스크립트 실행하기
\end{dhksmry}

\subsubsection{\zsh{subprocess.run()}과 \zsh{CompletedProcess}}

1. \zsh{subprocess.run()}은 무엇을 반환할까?

먼저 다음 코드를 살펴보자:

\begin{dhkcode}
import subprocess
import sys

result = subprocess.run(
    [sys.executable, "--version"]
)

print(result)
\end{dhkcode}

예를 들어 다음과 같은 결과가 나올 수 있다:

\begin{dhkzsh}
Python 3.12.7
CompletedProcess(args=['.../python', '--version'], returncode=0)
\end{dhkzsh}

여기서 두 가지가 출력되었다.

첫 번째는 자식 프로세스가 출력한 Python 버전이다.

\begin{dhkzsh}
Python 3.12.7
\end{dhkzsh}

두 번째는 \zsh{subprocess.run()}이 반환한 객체를 출력한 결과다.

\begin{dhkcode}
CompletedProcess(args=['.../python', '--version'], returncode=0)
\end{dhkcode}

즉, \zsh{subprocess.run()}은 단순히 외부 프로그램을 실행하는 것으로 끝나지 않는다.

외부 프로그램의 실행이 완료되면, 

실행 결과에 관한 정보를 담은 \zsh{CompletedProcess} 객체를 반환한다.

이 구조를 이전에 배운 함수의 관점에서 생각해 보자.

\begin{dhkcode}
result = subprocess.run(...)
\end{dhkcode}

여기서:

\begin{dhksmry}
    - subprocess.run 은 함수 객체를 참조한다.

    - subprocess.run(...)은 함수를 호출한다.

    - 함수가 반환한 CompletedProcess 객체가 
      result 라는 이름에 바인딩된다.
\end{dhksmry}

\zsh{Path} 객체나 \zsh{requests.Response} 객체와 비슷하게 이해하면 된다.

2. \zsh{CompletedProcess} 객체의 주요 속성

\zsh{CompletedProcess} 객체에는 대표적으로 다음 속성이 있다.

\begin{dhktable}
|속성|의미|
|:---|:---|
|args|실행한 프로그램과 인수|
|returncode|자식 프로세스의 종료 상태|
|stdout|자식 프로세스의 표준 출력 (캡쳐하도록 설정한 경우)|
|stderr|자식 프로세스의 표준 오류 출력 (캡쳐하도록 설정한 경우)|
\end{dhktable}

예제를 보자:

\begin{dhkcode}
import subprocess
import sys

result = subprocess.run(
    [sys.executable, "--version"],
    capture_output=True,
    text=True
)

print("args:", result.args)
print("returncode:", result.returncode)
print("stdout:", result.stdout)
print("stderr:", result.stderr)
\end{dhkcode}

예상 출력은 다음과 같다:

\begin{dhkzsh}
args: ['/path/to/python', '--version']
returncode: 0
stdout: Python 3.12.7
stderr:
\end{dhkzsh}

여기서 새로운 인수가 등장했다.

\begin{dhkcode}
(
    capture_output=True
    text=True
)
\end{dhkcode}

이 두 인수의 의미를 이해하려면 

먼저 운영체제의 \textbf{표준 입출력(Standard I/O)} 개념을 알아야 한다.

\subsubsection{표준 입력과 표준 출력}

프로세스는 외부와 데이터를 주고받을 수 있다.

그런데 프로세스마다 데이터를 입력받고 출력하는 방식을 완전히 다르게 

정의한다면, 프로그램을 연결하기가 매우 어려워진다.

그래서 운영체제는 일반적으로 다음과 같은 표준 입출력 통로를 제공한다.

\begin{dhktree}
Process
>stdin (Standard Input):: 표준 입력
>stdout (Standard Output):: 표준 출력
>stderr (Standard Error):: 표준 오류
\end{dhktree}

각각을 자세히 살펴본다.

\begin{dhknote}
\begin{dhksmry}[centered]
stdin - Standard Input
\end{dhksmry}
\end{dhknote}

프로세스가 데이터를 입력받는 통로이다.

가장 익숙한 예시는 \zsh{input()} 함수이다.

\begin{dhkcode}
name = input("이름을 입력하세요: ")
print(name)
\end{dhkcode}

터미널에서 실행하면 다음과 같이 동작한다.

\begin{dhkzsh}
이름을 입력하세요: Kim
Kim
\end{dhkzsh}

이때 일반적으로 키보드에서 입력한 데이터가 터미널을 통해 프로세스 \zsh{stdin}으로 전달된다.

다만 \zsh{stdin}이 반드시 키보드여야 하는 것은 아니다.

파일이나 다른 프로세스에서 데이터를 전달받을 수도 있다.

이것이 뒤에서 배울 파이프와 연결된다.

\begin{dhknote}
\begin{dhksmry}[centered]
stdout - Standard Output
\end{dhksmry}
\end{dhknote}

프로세스가 일반적인 실행 결과를 출력하는 통로이다.

\begin{dhkcode}
print("Training started")
\end{dhkcode}

우리가 사용하는 \zsh{print()}는 기본적으로 \zsh{sys.stdout}에 데이터를 출력한다.

\begin{dhkcode}
import sys

print("Hello")

sys.stdout.write("World\n")
\end{dhkcode}

결과:

\begin{dhkzsh}
Hello
World
\end{dhkzsh}

두 코드 모두 표준 출력으로 데이터를 보낸다는 점에서 비슷하다.

여기서 중요한 사실은 다음과 같다.

\zsh{stdout}은 터미널 화면 자체가 아니라 출력 데이터를 전달하는 통로이다.

보통 터미널에서 실행하면 \zsh{stdout}이 터미널에 연결되어 있어서 화면에 출력되는 것이다.

다른 프로그램이나 파일에 연결할 수도 있다.

\begin{dhknote}
\begin{dhksmry}[centered]
stderr - Standard Error
\end{dhksmry}
\end{dhknote}

오류 메시지나 진단 정보를 출력하는 별도의 통로이다.

예를 들어 다음 코드를 보자.

\begin{dhkcode}
import sys

print("Training started")

print(
    "Warning: GPU not found",
    file=sys.stderr
)
\end{dhkcode}

일반적인 터미널에서는 다음과 같이 출력될 수 있다.

\begin{dhkzsh}
Training started
Warning: GPU not found
\end{dhkzsh}

화면만 보면 둘 다 동일한 방식으로 출력된 것처럼 보인다.

하지만 내부적으로는 서로 다른 통로를 사용했다.

\begin{dhkdir}
Python Process
> stdout
>> Training started
> stderr
>> Warning: GPU not found
\end{dhkdir}

왜 굳이 분리할까?

예를 들어 ML 모델의 예측 결과를 다른 프로그램에 전달한다고 생각해 보자.

\begin{dhkzsh}
Prediction: 0.85
Warning: GPU not found
\end{dhkzsh}

두 종류의 출력이 섞이면 다른 프로그램이 예측 결과를 해석하기 어려울 수 있다.

그래서 정상적인 결과와 오류·진단 정보를 분리할 수 있도록 한 것이다.

단, 모든 프로그램이 이 관례를 엄격하게 지키는 것은 아니다.

\zsh{stderr}에 출력되었다고 반드시 실행 실패를 의미하지는 않는다.

\subsubsection{\zsh{capture_output=True}의 의미}

이제 다시 \zsh{subprocess.run()}으로 돌아가 보자.

다음 두 코드를 비교해 보자.

첫 번째 코드:

\begin{dhkcode}
import subprocess
import sys

result = subprocess.run(
    [sys.executable, "--version"]
)

print(result.stdout)
\end{dhkcode}

예상 결과:

\begin{dhkzsh}
Python 3.12.7
None
\end{dhkzsh}

왜 \zsh{None}이 출력될까?

자식 프로세스가 버전 정보를 출력하기는 했지만,

부모 Python 프로세스가 그 출력을 별도로 수집하지 않았기 때문이다.

기본적으로 자식 프로세스의 출력은 부모가 연결해 준 표준 출력 대상으로 전달된다.

따라서 터미널에서는 출력이 보이지만, \zsh{result.stdout}에 저장되는 것은 아니다.

두 번째 코드:

\begin{dhkcode}
import subprocess
import sys

result = subprocess.run(
    [sys.executable, "--version"],
    capture_output=True
)

print(result.stdout)
\end{dhkcode}

예상 결과:

\begin{dhkzsh}
b'Python 3.12.7\n'
\end{dhkzsh}

이번에는 \zsh{stdout}이 \zsh{None}이 아니라 데이터로 반환되었다.

왜냐하면:

\begin{dhkcode}
(
    capture_output=True
)
\end{dhkcode}

를 지정했기 때문이다.

이 옵션을 사용하면 \zsh{subprocess.run()}이

자식 프로세스의 \zsh{stdout}과 \zsh{stderr}를 캡처해 결과 객체에 저장한다.

그런데 출력에 \zsh{b}가 붙어 있다.

\begin{dhkzsh}
b'Python 3.12.7\n'
\end{dhkzsh}

이것은 bytes 객체이다.

운영체제의 표준 입출력에서 전달되는 데이터는 기본적으로 바이트 단위이기 때문에,

\zsh{subprocess}도 기본적으로 캡처 결과를 bytes로 반환한다.

따라서 다음과 같이 디코딩할 수 있다.

\begin{dhkcode}
print(result.stdout.decode("utf-8"))
\end{dhkcode}

하지만 매번 직접 디코딩하는 것은 번거롭다.

그래서 \zsh{text=True}를 제공한다.

\subsubsection{\zsh{text=True}와 문자열 변환}

다음 코드를 보자.

\begin{dhkcode}
import subprocess
import sys

result = subprocess.run(
    [sys.executable, "--version"],
    capture_output=True,
    text=True
)

print(type(result.stdout))
print(result.stdout)
\end{dhkcode}

결과:

\begin{dhkzsh}
<class 'str'>
Python 3.12.7
\end{dhkzsh}

\zsh{text=True}를 지정하면 \zsh{subprocess}가 입출력을 텍스트 모드로 처리한다.

캡처한 바이트를 문자열로 디코딩해서 반환한다.

\begin{dhktable}
|파일 입출력|subprocess|
|:---|:---|
|open("data.txt", "rb")|기본 캡처 모드|
|open("data.txt", "r", encoding='utf-8')|text=True, encoding='utf-8'|
|bytes 반환|기본적으로 bytes 반환|
|str 반환|text=True이면 str 반환|
\end{dhktable}

다음처럼 인코딩을 명시할 수도 있다.

\begin{dhkcode}
result = subprocess.run(
    [sys.executable, "--version"],
    capture_output=True,
    text=True,
    encoding="utf-8"
)
\end{dhkcode}

여기서 기억할 것은 단순하다.

\begin{dhknote}
\begin{dhksmry}
    - capture_output=True : 자식 프로세스의 출력을 수집한다.

    - text=True : 입출력 데이터를 문자열로 처리한다.
\end{dhksmry}
\end{dhknote}

두 옵션은 서로 다른 역할을 수행한다.

참고로 \zsh{capture_output=True} 없이 \zsh{text=True}만 지정해도

\zsh{result.stdout}이 자동으로 채워지는 것은 아니다.

\subsubsection{종료 코드(Exit Code)와 \zsh{returncode}}

이제 프로세스의 실행이 성공했는지 실패했는지 판단하는 방법을 알아보자.

외부 프로그램이 실행을 마치면 운영체제에 종료 상태를 전달할 수 있다.

이를 \zsh{종료 코드(Exit Code)}라고 부른다.

일반적으로 다음과 같은 관례를 사용한다.

\begin{dhktable}
|종료 코드|의미|
|:---|:---|
|0|정상 종료|
|0이 아닌 값|오류 또는 특정 종료 상태|
\end{dhktable}

여기서 주의할 점은 종료 코드의 구체적인 의미가 프로그램마다 다를 수 있다는 것이다.

예를 들어 어떤 프로그램은 \zsh{1}을 일반적인 오류에 사용하고,

다른 프로그램은 특정 조건을 나타내는 데 사용할 수도 있다.

\begin{dhknote}
\begin{dhksmry}[centered]
예제 1 - 정상 종료
\end{dhksmry}
\end{dhknote}

\begin{dhkcode}
import subprocess
import sys

result = subprocess.run(
    [sys.executable, "-c", "print('Hello')"],
    capture_output=True,
    text=True
)

print(result.stdout)
print(result.returncode)
\end{dhkcode}

결과:

\begin{dhkzsh}
Hello

0
\end{dhkzsh}

여기서 새로운 인수 \zsh{-c}가 등장했다.

\begin{dhkzsh}
python -c "print('Hello')"
\end{dhkzsh}

\zsh{-c}는 Python 인터프리터에게 뒤에 전달된 문자열을

Python 코드로 실행하라는 옵션이다.

즉, 별도의 \zsh{.py} 파일을 만들지 않고 간단한 Python 코드를 실행할 수 있다.

\begin{dhknote}
\begin{dhksmry}[centered]
예제 2 - 비정상 종료
\end{dhksmry}
\end{dhknote}

이번에는 의도적으로 오류를 발생시켜 보자.

\begin{dhkcode}
import subprocess
import sys

result = subprocess.run(
    [
        sys.executable,
        "-c",
        "raise ValueError('Invalid value')"
    ],
    capture_output=True,
    text=True
)

print("returncode:", result.returncode)
print("stdout:", result.stdout)
print("stderr:", result.stderr)
\end{dhkcode}

출력은 대략 다음과 같다:

\begin{dhkzsh}
returncode: 1
stdout:
stderr: Traceback (most recent call last):
...
ValueError: Invalid value
\end{dhkzsh}

여기서 중요한 부분이 있다.

자식 프로세스에서 \zsh{ValueError}가 발생했지만,

부모 Python 프로그램에서 \zsh{subprocess.run()}이 자동적으로 예외를 발생시키지 않았다.

대신 자식 프로세스의 종료 상태를 다음 속성에 저장했다.

\begin{dhkzsh}
result.returncode
\end{dhkzsh}

\begin{dhksmry}[centered]
즉, 자식 프로세스의 실패와 

부모 Python 프로세스의 예외 발생은 서로 다른 개념이다.
\end{dhksmry}

이 차이를 이해하는 것이 매우 중요하다.

\subsubsection{\zsh{check=True}와 \zsh{CalledProcessError}}

그렇다면 자식 프로세스가 실패했을 때 부모 프로그램에서도 예외가 발생하도록 할 수 있을까?

가능하다.

\begin{dhkcode}
import subprocess
import sys

result = subprocess.run(
    [
        sys.executable,
        "-c",
        "raise ValueError('Invalid value')"
    ],
    capture_output=True,
    text=True,
    check=True
)
\end{dhkcode}

여기서 추가된 옵션은 다음과 같다.

\begin{dhkcode}
(
    check=True
)
\end{dhkcode}

이 옵션을 지정하면 자식 프로세스가 0이 아닌 종료 코드를 반환했을 때

\zsh{subprocess.CalledProcessError} 예외가 발생한다.

따라서 다음과 같이 처리할 수 있다.

\begin{dhkcode}
import subprocess
import sys

try:
    result = subprocess.run(
        [
            sys.executable,
            "-c",
            "raise ValueError('Invalid value')"
        ],
        capture_output=True,
        text=True,
        check=True
    )

except subprocess.CalledProcessError as e:
    print("실행 실패")
    print("종료 코드:", e.returncode)
    print("오류 메시지:", e.stderr)
\end{dhkcode}

여기서 예외 객체 \zsh{e}는 자식 프로세스의 실패에 관한 정보를 가지고 있다.

실행 흐름을 정리하면 다음과 같다.

\begin{dhksmry}[centered]
Parent Python Process
\da
subprocess.run(check=True)
\da
Child Python Process
\da
ValueError 발생
\da
Child Process 종료
\da
Exit Code = 1
\da
Parent Python Process
\da
CalledProcessError 발생
\da
except 블록 실행
\end{dhksmry}

중요한 점은 부모에서 발생한 예외가 자식의 \zsh{ValueError} 객체 자체는 아니다.

부모가 받은 것은 

\zsh{자식 프로세스의 비정상 종료를 나타내는 CalledProcessError}다.

이 구조는 ML/DL 프로젝트에서 여러 학습 스크립트를 자동으로 실행할 때 특히 유용하다.

\subsubsection{\zsh{sys.exit()}과 프로세스 종료}

이제 프로세스의 종료 상태를 직접 지정하는 방법을 알아보자.

우리는 이미 \zsh{sys} 모듈을 배웠는데, \zsh{sys}에는 다음과 같은 함수도 있다.

\begin{dhkcode}
sys.exit()
\end{dhkcode}

이 함수는 현재 Python 프로그램의 실행을 종료하도록 요청하는 역할을 한다.

예를 들어 다음 코드를 보자:

\begin{dhkcode}
# train.py

import sys

print("Training started")

sys.exit(0)

print("Training finished")
\end{dhkcode}

실행 결과:

\begin{dhkzsh}
Training started
\end{dhkzsh}

\zsh{sys.exit(0)} 이후의 코드는 실행되지 않는다.

여기서 \zsh{0}은 정상 종료를 나타내는 종료 코드로 사용된다.

이번에는 다음처럼 작성해 보자:

\begin{dhkcode}
# train.py

import sys

print("Training started")

sys.exit(1)
\end{dhkcode}

이 프로그램을 \zsh{subprocess.run()}으로 실행하면 다음과 같다:

\begin{dhkcode}
import subprocess
import sys

result = subprocess.run(
    [sys.executable, "train.py"]
)

print(result.returncode)
\end{dhkcode}

결과:

\begin{dhkzsh}
Training started
1
\end{dhkzsh}

\zsh{train.py}가 \zsh{sys.exit(1)}로 종료되었기 때문에

부모 프로그램은 \zsh{returncode}에서 \zsh{1}을 확인할 수 있다.

\begin{dhknote}
\begin{dhksmry}[centered]
sys.exit()은 내부적으로 무엇을 할까?
\end{dhksmry}
\end{dhknote}

사실 \zsh{sys.exit()}은 내부적으로 \zsh{SystemExit} 예외를 발생시킨다.

\begin{dhkcode}
sys.exit(1)
\end{dhkcode}

는 개념적으로 다음과 유사하다:

\begin{dhkcode}
raise SystemExit(1)
\end{dhkcode}

따라서 일반적인 예외처럼 처리될 여지가 있고, \zsh{finally} 블록도 실행된다.

\begin{dhkcode}
import sys

try:
    print("A")
    sys.exit(1)

finally:
    print("B")
\end{dhkcode}

결과:

\begin{dhkzsh}
A
B
\end{dhkzsh}

그 후 \zsh{SystemExit}이 처리되지 않은 채 최상위까지 전달되면 Python 인터프리터가 종료된다.

즉, \zsh{sys.exit()}은 운영체제에 프로세스를 즉시 강제 종료하라는 명령과는 다르다.

이 차이 때문에 일반적인 Python 프로그램에서는 \zsh{sys.exit()}을 사용하고,

\zsh{os._exit()}은 특별한 상황에서만 사용한다.

\zsh{os._exit()}은 Python의 일반적인 정리 절차를 거치지 않고 프로세스를 

즉시 종료하기 때문에 보통의 학습, 응용 코드에서 사용할 필요는 없다.

\subsubsection{실제 ML 프로젝트에서 활용하기}

이제 지금까지 배운 내용을 실제 프로젝트에 연결해 보자.

다음과 같은 프로젝트가 있다고 가정하자.

\begin{dhkdir}
ml-project/
>train.py
>run_experiment.py
>data/
>>train.csv
>results/
\end{dhkdir}

\zsh{train.py}는 학습을 수행하는 스크립트다.

\begin{dhkcode}
# train.py

import argparse
import sys
from pathlib import Path

parser = argparse.ArgumentParser()

parser.add_argument(
    "--data",
    type=Path,
    required=True
)

parser.add_argument(
    "--epochs",
    type=int,
    default=10
)

args = parser.parse_args()

if not args.data.is_file():
    print(
        f"Data file not found: {args.data}",
        file=sys.stderr
    )

    sys.exit(1)

print(f"Data: {args.data}")
print(f"Epochs: {args.epochs}")

print("Training completed")
\end{dhkcode}

이 코드는 지금까지 배운 개념들로 구성되어 있다.

\begin{dhksmry}
    argparse    : 커맨드 라인 인수 처리

    Path        : 파일 경로 표현

    is_file()   : 파일 존재 여부 확인

    sys.stderr  : 오류 메시지 출력

    sys.exit(1) : 오류 상태로 종료
\end{dhksmry}

이번에는 \zsh{run_experiment.py}를 작성해 보자:

\begin{dhkcode}
# run_experiment.py

import subprocess
import sys
from pathlib import Path

project_root = Path(__file__).resolve().parent

train_script = project_root / "train.py"
data_path = project_root / "data" / "train.csv"

command = [
    sys.executable,
    str(train_script),
    "--data",
    str(data_path),
    "--epochs",
    "10"
]

result = subprocess.run(
    command,
    capture_output=True,
    text=True,
    encoding="utf-8"
)

print("Exit code:", result.returncode)
print("Standard output:")
print(result.stdout)

print("Standard error:")
print(result.stderr)
\end{dhkcode}

여기서 눈여겨볼 부분은 다음이다.

\begin{dhkcode}
command = [
    sys.executable,
    str(train_script),
    "--data",
    str(data_path),
    "--epochs",
    "10"
]
\end{dhkcode}

이전에 배운 커맨드라인 인수가 등장했다.

\zsh{subprocess.run()}의 리스트는 실행 프로그램과 각 인수를 구분해서 전달하는 역할을 한다.

또한 \zsh{"10"}이 문자열인 이유도 이해할 수 있다.

커맨드라인 인수는 문자열로 전달되고, \zsh{train.py}의 \zsh{argparse}가 다음 설정을 통해 정수로 변환한다.

\begin{dhkcode}
parser.add_argument(
    "--epochs",
    type=int
)
\end{dhkcode}

전체 흐름을 보면 다음과 같다.

\begin{dhkbranch}[in=8]
node main :: run_experiment.py
arrow main :: subprocess.run()
node main :: Child Python Process
arrow main
node main :: train.py
branchline main -> argparse :: argparse //* --data //* --epochs
stop argparse
branchline main -> path :: Path //* 데이터 파일 검증
stop path
branchline main -> stdout :: stdout //* 학습 결과 출력
stop stdout
branchline main -> stderr :: stderr //* 오류 메시지 출력
stop stderr
handoff main -> exit :: Exit code
arrow exit
node exit :: run_experiment.py //* result.stdout //* result.stderr //* result.returncode
\end{dhkbranch}

이것이 Python으로 핛흡 실행을 자동화하는 가장 기본적인 구조이다.

물론 실제 모델 학습에서는 \zsh{subprocess.run()} 대신 

하나의 Python 프로세스 안에서 학습 함수를 직접 호출하는 방식도 사용할 수 있다.

별도 프로세스를 실행하는 것이 항상 더 좋은 것은 아니고,

실행 환경 격리나 독립적인 스크립트 실행이 필요할 때 유용하다.

\subsubsection{이번 학습 내용 정리}

지금까지 살펴본 내용을 표로 정리해 보자.

\begin{dhktable}
|개념|역할|
|:---|:---|
|subprocess.run()|외부 프로그램을 실행하고 종료를 기다림|
|CompletedProcess|실행 결과 정보를 담은 객체|
|args|실행한 명령과 인수|
|stdin|표준 입력|
|stdout|표준 출력|
|stderr|표준 오류 출력|
|capture_output=True|자식 프로세스의 출력을 캡처|
|text=True|입출력을 문자열로 처리|
|returncode|자식 프로세스의 종료 상태|
|check=True|비정상 종료 시 CalledProcessError 발생|
|sys.exit()|SystemExit를 발생시켜 프로그램 종료 요청|
|sys.executable|현재 Python 인터프리터의 실행 파일 경로|
\end{dhktable}

\begin{dhknote}
\begin{dhksmry}[centered]
이번 장에서 가장 중요한 세 가지
\end{dhksmry}
\end{dhknote}

첫 번째, 프로그램의 출력과 함수의 반환값은 서로 다르다

\begin{dhkcode}
result = subprocess.run(...)
\end{dhkcode}

여기서 자식 프로세스가 출력하는 내용과

\zsh{subprocess.run()}이 반환하는 \zsh{CompletedProcess} 객체는 별개의 개념이다.

두 번째, 자식 프로세스의 실패와 부모 프로세스의 예외는 서로 다르다.

\begin{dhkcode}
result.returncode
\end{dhkcode}

는 자식 프로세스의 종료 상태를 나타내고,

\begin{dhkcode}
(
    check=True
)
\end{dhkcode}

는 그 상태를 부모 Python 프로그램의 예외 처리와 연결해주는 옵션이다.

세 번째, 표준 입출력은 터미널 화면이 아니라 데이터 통로다.

\begin{dhkdiagram}
node stdin (1,1) :: stdin
node process (1,2) :: Process
node stdout (1,3) :: stdout
node stderr (2,3) :: stderr

edge stdin -> process
edge process -> stdout
edge process -> stderr :: VH
\end{dhkdiagram}

이 구조를 이해해야 서로 다른 프로그램을 연결하는 파이프와 리다이렉션을 이해할 수 있다.

다음 학습에서는 Shell의 파이프와 리다이렉션,

다음과 같은 터미널에서 사용하는 기호들의 의미를 살펴본다.

\begin{dhkzsh}
python train.py > train.log
python train.py 2> error.log
python train.py > train.log 2>&1
python generate.py | python analyze.py
\end{dhkzsh}

여기서 \zsh{>, 2>, 2>&1, |}는 

Python 문법이 아니라 Shell이 제공하는 기능이다.

이 부분까지 이해하면 Shell, Python, 운영체제, 프로세스, 파일 입출력이 어떻게 연결되는지 거의 완성된다.

\subsubsection{Shell과 표준 입출력의 관계}

먼저 지금가지 배운 프로세스 구조를 다시 살펴보자

\begin{dhkbranch}[in=8]
node main :: Operating System
arrow main
node main :: Shell Process
arrow main :: python train.py
node main :: Python Process
branchline main -> in :: stdin
stop in
branchline main -> out :: stdout
stop out
handoff main -> err :: stderr
\end{dhkbranch}

여기서 중요한 질문이 하나 있다.

Python 프로세스가 \zsh{print()}를 실행하면 왜 터미널 화면에 문자열이 출력될까?

우리는 지금까지 다음과 같이 이해했다.

\begin{dhkcode}
print("Hello")
\end{dhkcode}

\zsh{print()}는 기본적으로 \zsh{sys.stdout}에 문자열을 출력한다.

그런데 \zsh{stdout} 자체가 터미널 화면을 의미하는 것은 아니라고 하였다.

그렇다면 다음 두 개념을 구분해야 한다.

\begin{dhksmry}
    1. 표준 출력(stdout)
        프로세스가 데이터를 내보내는 논리적 통로

    2. 터미널(Terminal)
        해당 데이터를 화면에 표시할 수 있는 입출력 환경
\end{dhksmry}

일반적인 터미널 실행에서는 Shell이 자식 프로세스를 생성할 때 표준 입출력이 터미널에 연결되도록 한다.

\begin{dhkdiagram}
node ter (1,4) :: Terminal
node std (2,2) :: stdout
node pro (1,1) :: Python Process

edge pro -> ter
\end{dhkdiagram}

따라서 Python이 \zsh{stdout}에 데이터를 쓰면 터미널에 표시된다.

하지만 Shell은 이 연결을 변경할 수도 있다.

예를 들어 \zsh{stdout}을 터미널이 아니라 파일에 연결할 수 있다.

이것을 \zsh{리다이렉션(Redirection)}이라고 한다.

\subsubsection{파일 디스크립터(File Descriptor)}

리다이렉션을 이해하기 전에 새로운 개념 하나를 알아야 한다.

바로 \zsh{파일 디스크립터(File Descriptor, FD)}이다.

운영체제는 프로세스가 열어 놓은 파일이나 입출력 자원에 접근할 수 있도록 식별자를 제공한다.

POSIX 계열 운영체제(macOS/Linux)에서는 이러한 식별자로 정수를 사용한다.

대표적인 표준 파일 디스크립터는 다음과 같다.
\begin{dhktable}
|파일 디스크립터	|이름|	의미|
|:---|:---|:---|
|0|stdin	|표준 입력|
|1|stdout	|표준 출력|
|2|stderr|표준 오류 출력|
\end{dhktable}

이를 구조로 표현하면 다음과 같다.

\begin{dhkdir}
Python Process
> FD 0 -> stdin
> FD 1 -> stdout
> FD 2 -> stderr
\end{dhkdir}

여기서 한 가지 주의해야 한다.

\zsh{0, 1, 2}는 Python에서 정의한 숫자가 아니다.

POSIX 계열 운영체제에서 사용하는 표준 파일 디스크립터 번호다.

Python은 이러한 운영체제 기능을 활용하는 것이지,

Windows에서도 표준 입출력 스트림이라는 개념은 존재하지만,

내부적인 핸들 관리 방식은 POSIX와 차이가 있다.

\begin{dhknote}
\begin{dhksmry}[centered]
Python에서 확인하기
\end{dhksmry}
\end{dhknote}

macOS의 일반적인 터미널 환경에서 다음 코드를 실행해 보자.

\begin{dhkcode}
import sys

print(sys.stdin.fileno())
print(sys.stdout.fileno())
print(sys.stderr.fileno())
\end{dhkcode}

일반적으로 다음 결과가 나온다.

\begin{dhkzsh}
0
1
2
\end{dhkzsh}

\zsh{fileno()}는 해당 스트림과 연결된 운영체제 수준의 파일 디스크립터 번호를 반환하는 메서드다.

다만 Jupyter Notebook이나 일부 IDE에서는 표준 스트림을 별도의 객체로 대체하기 때문에 \zsh{fileno()}가 지원되지 않을 수 있다.

여기서 핵심은 다음이다.

\begin{dhksmry}[centered]
Shell은 파일 디스크립터의 연결 대상을 변경하여

프로세스의 입출력 방향을 제어할 수 있다.
\end{dhksmry}

\subsubsection{출력 리다이렉션 \zsh{>}}

이제 가장 기본적인 리다이렉션부터 살펴보자.

다음 Python 파일이 있다고 가정하자.

\begin{dhkcode}
# train.py

print("Training started")
print("Epoch 1 completed")
print("Epoch 2 completed")
\end{dhkcode}

터미널에서 다음과 같이 실행하면:

\begin{dhkzsh}
python train.py
\end{dhkzsh}

결과가 화면에 출력된다:

\begin{dhkzsh}
Training started
Epoch 1 completed
Epoch 2 completed
\end{dhkzsh}

그런데 다음처럼 실행하면 어떻게 될까?

\begin{dhkzsh}
python train.py > train.log
\end{dhkzsh}

이번에는 정상 출력이 터미널에 표시되지 않고 \zsh{train.log} 파일에 저장된다.

파일 내용은 다음과 같다:

\begin{dhkzsh}
Training started
Epoch 1 completed
Epoch 2 completed
\end{dhkzsh}

이것이 출력 리다이렉션이다.

\begin{dhknote}
\begin{dhksmry}[centered]
내부 구조
\end{dhksmry}
\end{dhknote}

일반적인 실행:

\begin{dhkdiagram}[in=6]
node pp (1,1) :: Python Process
node fd1 (2,2) :: FD 1 (stdout)
node tm (3,2) :: Terminal

edge pp -- fd1 :: VH
edge fd1 -> tm
\end{dhkdiagram}

리다이렉션을 사용한 실행:

\begin{dhkdiagram}[in=6]
node pp (1,1) :: Python Process
node fd1 (2,2) :: FD 1 (stdout)
node tm (3,2) :: train.log

edge pp -- fd1 :: VH
edge fd1 -> tm
\end{dhkdiagram}

여기서 중요한 사실이 있다.

\begin{dhksmry}[centered]
Python 코드 자체는 변경되지 않았다.
\end{dhksmry}

\begin{dhkcode}
print("Training started")
\end{dhkcode}

이 코드는 여전히 표준 출력에 데이터를 보내고 있다.

다만 Shell이 프로세스를 실행할 때 \zsh{stdout}의 연결 대상을 파일로 변경했기 때문에 파일에 저장되는 것이다.

즉, 출력 목적지를 결정하는 주체가 Python 코드가 아니라 Shell인 경우이다.

\begin{dhknote}
\begin{dhksmry}[centered]
\zsh{>}는 기존 파일을 어떻게 처리할까?
\end{dhksmry}
\end{dhknote}

다음 명령을 두 번 실행해 보자.

\begin{dhkzsh}
python train.py > train.log
\end{dhkzsh}

\begin{dhkzsh}
python train.py > train.log
\end{dhkzsh}

일반적으로 두 번째 명령을 실행하면 기존 \zsh{train.log}의 내용이 지워지고 새로운 내용으로 대체된다.

왜냐하면 \zsh{>}는 기본적으로 파일을 생성하거나 기존 파일을 잘라내고(truncate) 출력을 기록하기 때문이다.

이전에 배운 다음 코드와 비슷한 성질이다.

\begin{dhkcode}
with open("train.log", "w") as f:
    f.write("Training started\n")
\end{dhkcode}

다만 Shell 리다이렉션과 Python의 \zsh{open()}은 서로 다른 계층에서 동작한다는 차이가 있다.

\subsubsection{추가 출력 리다이렉션 \zsh{>>}}

기존 파일 내용을 유지하면서 새로운 내용을 추가하고 싶다면 어떻게 해야 할까?

이때는 \zsh{>>}를 사용한다.

\begin{dhkzsh}
python train.py >> train.log
\end{dhkzsh}

두 연산자를 비교하면 다음과 같다.

\begin{dhktable}
|연산자	|동작	|Python 파일 모드와 비교|
|:---|:---|:---|
|>|	파일을 새로 쓰거나 기존 내용을 잘라내고 기록	|"w"|
|> >|	기존 파일 끝에 내용을 추가	|"a"|
\end{dhktable}

예를 들어 다음과 같이 실행한다고 생각해 보자.

\begin{dhkzsh}
python train.py > train.log
python train.py >> train.log
\end{dhkzsh}

그러면 \zsh{train.log}에는 두 번의 실행 결과가 순서대로 기록된다.

\begin{dhkzsh}
Training started
Epoch 1 completed
Epoch 2 completed
Training started
Epoch 1 completed
Epoch 2 completed
\end{dhkzsh}

ML/DL 프로젝트에서 학습 로그를 누적할 때 사용할 수 있다.

하지만 여러 실험을 하나의 로그 파일에 계속 추가하면 

어느 실험의 결과인지 구분하기 어려워질 수 있으므로, 

실제 프로젝트에서는 실행별 로그 파일을 분리하는 방법도 자주 사용한다.

\subsubsection{오류 출력 리다이렉션 \zsh{2>}}

이번에는 \zsh{stderr}를 살펴보자

다음 코드를 작성했다고 가정하자:

\begin{dhkcode}
# train.py

import sys

print("Training started")

print(
    "Error: Dataset not found",
    file=sys.stderr
)
\end{dhkcode}

일반적으로 실행하면 두 메시지가 모두 터미널에 표시된다.

\begin{dhkzsh}
Training started
Error: Dataset not found
\end{dhkzsh}

하지만 내부적으로는 서로 다른 통로를 사용하고 있다.

\begin{dhkdir}
Python Process
> FD 1 (stdout)
>> Training started
> FD 2 (stderr)
>> Error: Dataset not found
\end{dhkdir}

그렇다면 다음 명령을 실행해 보자.

\begin{dhkzsh}
python train.py > train.log
\end{dhkzsh}

어떻게 될까?

\zsh{stdout}만 파일로 연결했기 때문에 \zsh{train.log}에는 다음 내용만 저장된다.

\begin{dhkzsh}
Training started
\end{dhkzsh}

반면 오류 메시지는 여전히 터미널에 표시된다.

\begin{dhkzsh}
Error: Dataset not found
\end{dhkzsh}

이번에는 다음 명령을 실행해 보자.

\begin{dhkzsh}
python train.py 2> error.log
\end{dhkzsh}

여기서 \zsh{2}는 \zsh{stderr}의 파일 디스크립터 번호다.

따라서 다음과 같이 동작한다.

\begin{dhkdiagram}[in=6]
node pp (1,1) :: Python Process
node out (2,2) :: FD 1 (stdout)
node ter (3,2) :: Terminal
node err (4,2) :: FD 2 (stderr)
node log (5,2) :: error.log

edge pp -- out :: VH
edge out -> ter
edge pp -- err :: VH
edge err -> log
\end{dhkdiagram}

\zsh{stdout}은 터미널에 표시되고, \zsh{stderr}는 \zsh{error.log}에 저장된다.

\begin{dhknote}
\begin{dhksmry}[centered]
두 출력을 서로 다른 파일에 저장하기
\end{dhksmry}
\end{dhknote}

다음처럼 작성할 수도 있다.

\begin{dhkzsh}
python train.py > train.log 2> error.log
\end{dhkzsh}

그러면 다음과 같이 분리된다.

\begin{dhkdir}
Python Process
> stdout -> train.log
> stderr -> error.log
\end{dhkdir}

이 방식은 프로그램의 정상 실행 결과와 오류 메시지를 구분해 저장하고 싶을 때 유용하다.

\subsubsection{\zsh{2>&1}의 의미}

이제 처음보면 이해하기 어려운 표현이 등장한다

\begin{dhkzsh}
python train.py > train.log 2>&1
\end{dhkzsh}

이것을 단순히 암기하기 보다는 기호를 나누어서 이해해 보자.

\begin{dhkdiagram}
node oo (1,1) :: 2
node otw (1,2) :: >
node ot (1,3) :: &
node of (1,4) :: 1
node f1 (2,5) :: 파일 디스크립터 1
node f2 (3,5) :: 파일 디스크립터를 참조
node f3 (4,5) :: 리다이렉션
node f4 (5,5) :: 파일 디스크립터 2

edge of -- f1 :: VH
edge ot -- f2 :: VH
edge otw -- f3 :: VH
edge oo -- f4 :: VH 
\end{dhkdiagram}

즉, \zsh{2>&1}은 FD 2를 현재 FD 1이 가리키는 출력 대상으로 연결하라는 의미다.

여기서 중요한 단어는 '현재'다.

\begin{dhknote}
\begin{dhksmry}[centered]
실행 순서
\end{dhksmry}
\end{dhknote}

다음 명령을 보자.

\begin{dhkzsh}
python trian.py > train.log 2>&1
\end{dhkzsh}

Shell은 리다이렉션을 왼쪽에서 오른쪽 순서로 처리한다.

먼저:

\begin{dhkzsh}
> train.log
\end{dhkzsh}

를 처리하면 다음과 같다.

\begin{dhksmry}
FD 1 \ra train.log
FD 2 \ra Terminal
\end{dhksmry}

그다음:

\begin{dhkzsh}
2>&1
\end{dhkzsh}

을 처리하면 FD 2도 현재 FD 1이 연결된 대상으로 향하게 된다.

\begin{dhkdiagram}
node fd1 (1,1) :: FD 1
node fd2 (3,1) :: FD 2
node log (2,2) :: train.log

edge fd1 -> log :: HV
edge fd2 -> log :: HV
\end{dhkdiagram}

따라서 정상 출력과 오류 출력이 모두 \zsh{train.log}에 저장된다.

\begin{dhknote}
\begin{dhksmry}[centered]
순서를 바꾸면 어떻게 될까?
\end{dhksmry}
\end{dhknote}

이번에는 다음 명령을 살펴보자.

\begin{dhkzsh}
python train.py 2>&1 > train.log
\end{dhkzsh}

비슷해 보이지만 결과가 다르다.

첫 번째 단계:

\begin{dhkzsh}
2>&1
\end{dhkzsh}

이 시점에서 FD 1은 터미널에 연결되어 있다.

따라서 FD 2도 터미널에 연결된다.

\begin{dhksmry}
FD 1 \ra Terminal
FD 2 \ra Terminal
\end{dhksmry}

두 번째 단계:

이제 FD 1만 파일에 연결된다.

\begin{dhkzsh}
FD 1 \ra train.log
FD 2 \ra Terminal
\end{dhkzsh}

따라서 최종 결과는 다음과 같다.

\begin{dhktable}
|명령|	stdout|	stderr|
|:---|:---|:---|
|> train.log 2>&1	|파일|	파일|
|2>&1 > train.log	|파일	|터미널|
\end{dhktable}

이것이 Shell 리다이렉션에서 연산자 순서가 중요한 이유다.

특히 \zsh{2>&1}은 두 스트림을 영구적으로 하나의 스트림으로 합친다는 뜻이 아니라, FD 2의 출력 대상을 그 시점의 FD 1과 동일하게 설정한다는 의미로 이해하는 것이 정확하다.

\subsubsection{입력 리다이렉션 \zsh{<}}

지금까지는 프로세스의 출력을 다른 곳으로 보냈다.

이번에는 반대로 프로세스에 데이터를 입력하는 방법을 알아보자.

다음 파일을 작성해 보자.

\begin{dhkcode}
# analyze.py

import sys

data = sys.stdin.read()

print("Received:")
print(data)
\end{dhkcode}

이 프로그램은 \zsh{stdin}으로 전달된 데이터를 모두 읽는다.

일반적으로 터미널에서 실행하면 입력을 기다리게 된다.

\begin{dhkzsh}
python analyze.py
\end{dhkzsh}

이때 키보드로 문자열을 입력할 수 있다.

그런데 다음과 같이 실행하면 어떨까?

\begin{dhkzsh}
python analyze.py < data.txt
\end{dhkzsh}

이 명령은 \zsh{data.txt}를 표준 입력에 연결한다.

\begin{dhksmry}[centered]
data.txt
\da
FD 0 (stdin)
\da
Python Process
\da
sys.stdin.read()
\end{dhksmry}

따라서 Python 코드를 수정하지 않아도 파일의 내용을 입력으로 사용할 수 있다.

\begin{dhknote}
\begin{dhksmry}[centered]
Path.read_text()와 무엇이 다를까?
\end{dhksmry}
\end{dhknote}

다음 두 코드를 비교해 보자.

방법 A - 파일 경로를 직접 사용

\begin{dhkcode}
from pathlib import Path

data = Path("data.txt").read_text()
\end{dhkcode}

이 경우 Python 코드가 특정 파일을 직접 연다.

방법 B - 표준 입력 사용

\begin{dhkcode}
import sys

data = sys.stdin.read()
\end{dhkcode}

이 경우 Python은 입력 데이터가 어디에서 오는지 알 필요가 없다.

Shell이 파일을 연결할 수도 있고, 사용자가 키보드로 입력할 수도 있고, 다른 프로세스가 데이터를 전달할 수도 있다.

이것이 표준 입력을 사용하는 중요한 장점이다.

입력 데이터의 출처와 데이터를 처리하는 프로그램을 분리할 수 있다는 것이다.

\subsubsection{파이프(Pipe) \zsh{|}}

이제 이번 학습에서 가장 중요한 개념으로 넘어간다.

다음 명령을 살펴보자.

\begin{dhkzsh}
python generate.py | python analyze.py
\end{dhkzsh}

여기서 \zsh{|}는 Shell의 파이프 연산자다.

파이프는 일반적으로 앞 프로세스의 표준 출력을 뒤 프로세스의 표준 입력에 연결한다.

구조는 다음과 같다:

\begin{dhkbranch}
node main :: Python Process A
node main :: generate.py
arrow main :: stdout
node main :: PIPE
arrow main :: stdin
node main :: Python Process B
node main :: analyze.py
\end{dhkbranch}

여기서 중요한 점은 파일을 중간에 저장하지 않아도 된다는 것이다.

프로세스 A가 출력한 데이터를 프로세스 B가 입력받을 수 있다.

예제 1. 데이터 생성

먼저 \zsh{generate.py}를 작성해 보자:

\begin{dhkcode}
# generate.py

for i in range(1, 6):
    print(i)
\end{dhkcode}

실행하면 다음과 같이 출력된다.

\begin{dhkzsh}
1
2
3
4
5
\end{dhkzsh}

예제 2. 데이터 분석

이번에는 \zsh{analyze.py}를 작성해 보자:

\begin{dhkcode}
# analyze.py

import sys

numbers = []

for line in sys.stdin:
    numbers.append(int(line.strip()))

print("Sum:", sum(numbers))
\end{dhkcode}

이 코드도 새로운 내용은 거의 없다.

파일 객체의 반복 처리와 동일한 방식으로 \zsh{sys.stdin}을 사용하고 있다.

\begin{dhkcode}
for line in sys.stdin:
\end{dhkcode}

\zsh{sys.stdin}은 파일과 유사한 인터페이스를 제공하기 때문에 줄 단위 반복이 가능하다.

그리고 다음 코드로 개행 문자 등을 제거한다.

\begin{dhkcode}
line.strip()
\end{dhkcode}

이제 두 프로그램을 연결해 보자:

\begin{dhkzsh}
python generate.py | python analyze.py
\end{dhkzsh}

출력:

\begin{dhkzsh}
Sum: 15
\end{dhkzsh}

전체 흐름을 표현하면 다음과 같다:

\begin{dhkbranch}[in=8]
node main :: generate.py
line main :: print(1)
line main :: print(2)
line main :: print(3)
line main :: print(4)
arrow main :: print(5)
node main :: stdout
arrow main
node main :: PIPE
arrow main
node main :: stdin
arrow main
node main :: analyze.py
line main :: for line in sys.stdin
line main :: int(line.strip())
arrow main :: sum(numbers)
node main :: Sum: 15
\end{dhkbranch}

이 구조에서 두 프로그램은 서로의 내부 구현을 알 필요가 없다.

\zsh{generate.py}는 데이터를 출력하기만 하면 되고,

\zsh{analyze.py}는 입력된 데이터를 처리하기만 하면 된다.

즉, 표준 입출력이라는 공통 인터페이스를 통해 서로 다른 프로그램을 조합할 수 있는 것이다.

\subsubsection{파이프와 함수 호출의 차이}

여기서 Python의 함수 호출과 비교해 보자.

우리가 일반적으로 작성하는 코드는 다음과 같다.

\begin{dhkcode}
def generate():
    return [1, 2, 3, 4, 5]


def analyze(numbers):
    return sum(numbers)


data = generate()
result = analyze(data)

print(result)
\end{dhkcode}

이 경우에는 하나의 Python 프로세스 안에서 함수들이 데이터를 전달한다.

\begin{dhkdiagram}[in=6]
node pp (1,1) :: Python Process
node gen (2,2) :: generate()
node list (3,2) :: list 객체
node an (4,2) :: analyze()

edge pp -- gen :: VH
edge pp -- an :: VH
edge gen -> list
edge list -> an
\end{dhkdiagram}

반면 파이프는 서로 다른 프로세스 사이에서 데이터를 전달한다.

\begin{dhkbranch}
node main :: Process A
arrow main :: bytes
node main :: PIPE
arrow main :: bytes
node main :: Process B
\end{dhkbranch}

여기에 핵심적인 차이가 있다.

함수 호출은 같은 프로세스 안에서 Python 객체를 전달할 수 있지만,

일반적인 운영체제 파이프는 바이트 데이터를 전달한다.

따라서 파이프를 통해 리스트나 딕셔너리 같은 Python 객체 자체가 그대로 전달되는 것은 아니다.

예를 들어 다음 코드를 보자:

\begin{dhkcode}
print([1, 2, 3])
\end{dhkcode}

이 코드는 리스트 객체를 \zsh{stdout}으로 전달하는 것이 아니다.

리스트의 문자열 표현을 출력하는 것이다.

파이프를 통해 전달되는 것은 그 문자열을 인코딩한 바이트 데이터다.

그래서 실제 데이터 처리 파이프라인에서는 JSON, CSV< JSON Lines 같은 직렬화 형식을 사용하기도 한다.

\subsubsection{Shell 파이프와 Python \zsh{subprocess.PIPE}}

여기까지 이해했다면 이전에 배운 \zsh{subprocess.run()}과 연결할 수 있다.

다음 Shell 명령을 생각해 보자:

\begin{dhkzsh}
python generate.py | python analyze.py
\end{dhkzsh}

Python에서도 유사한 작업을 구현할 수 있다.

먼저 단순한 방법을 보자:

\begin{dhkcode}
import subprocess
import sys

generated = subprocess.run(
    [sys.executable, "generate.py"],
    capture_output=True,
    text=True,
    check=True
)

analyzed = subprocess.run(
    [sys.executable, "analyze.py"],
    input=generated.stdout,
    capture_output=True,
    text=True,
    check=True
)

print(analyzed.stdout)
\end{dhkcode}

결과:

\begin{dhkzsh}
Sum: 15
\end{dhkzsh}

여기서 새로운 인수는 다음과 같다:
\begin{dhkcode}
input = generated.stdout
\end{dhkcode}

이 인수는 자식 프로세스의 표준 입력으로 데이터를 전달한다.

실행 흐름은 다음과 같다:
\begin{dhksmry}[centered]
Parent Python Process
\da
generate.py 실행
\da
stdout 캡처
\da
generated.stdout
\da
analyze.py 실행
\da
stdin으로 데이터 전달
\da
stdout 캡처
\da
Parent Python Process
\end{dhksmry}

하지만 여기에는 Shell 파이프와 차이가 있다.

이 코드는 첫 번째 프로세스가 완전히 종료된 다음, 캡처한 출력을 두 번째 프로세스에 전달한다.

반면 Shell 파이프는 일반적으로 두 프로세스를 동시에 실행하고, 한쪽에서 생성한 데이터를 다른쪽에서 읽을 수 있도록 연결한다.

따라서 위의 \zsh{subprocess.run()} 두 번은 파이프와 비슷한 데이터 전달을 구현한 것이지만, 실제로 두 프로세스를 스트리밍 방식으로 연결한 것은 아니다.

실제 운영체제 파이프를 직접 연결하려면 \zsh{subprocess.Popen()}과 

\zsh{subprocess.PIPE} 등을 사용할 수 있다.

\subsubsection{\zsh{|}와 \zsh{>}를 함께 사용하기}

파이프와 리다이렉션을 결합할 수도 있다.

\begin{dhkzsh}
python generate.py | python analyze.py > result.txt
\end{dhkzsh}

이 명령을 해석해 보자.

첫 번째 단계:

\begin{dhkzsh}
python generate.py
\end{dhkzsh}

숫자 데이터를 생성한다.

두 번째 단계:
\begin{dhkzsh}
| python analyze.py
\end{dhkzsh}

생성된 데이터를 분석 프로그램의 표준 입력으로 전달한다.

세 번째 단계:

\begin{dhkzsh}
> result.txt
\end{dhkzsh}

분석 프로그램의 표준 출력을 파일에 저장한다.

전체 구조는 다음과 같다:

\begin{dhkbranch}[in=8]
node main :: generate.py
arrow main :: stdout
node main :: PIPE
arrow main :: stdin
node main :: analyze.py
arrow main :: stdout
node main :: result.txt
\end{dhkbranch}

결과적으로 \zsh{result.txt}에는 다음 내용이 저장된다.

\begin{dhkzsh}
Sum: 15
\end{dhkzsh}

이런 방식으로 여러 프로그램을 조합해 하나의 데이터 처리 흐름을 구성할 수 있다.

\subsubsection{파이프와 종료 코드의 관계}

마지막으로 주의할 부분이 하나 있다.

다음 명령을 보자.

\begin{dhkzsh}
python generate.py | python analyze.py
\end{dhkzsh}

여기에는 두 개의 프로세스가 존재한다.

따라서 각 프로세스는 독립적인 종료 코드를 가질 수 있다.

\begin{dhkdir}
generate.py
> Exit Code = 1
\end{dhkdir}

\begin{dhkdir}
analyze.py
> Exit Code = 0
\end{dhkdir}

그런데 Shell에서 파이프라인 전체의 종료 코드를 확인하면 어떻게 될까?

일반적인 Bash와 zsh에서는 기본적으로 마지막 명령의 종료 코드가 파이프라인의 종료 상태가 된다.

따라서 앞의 \zsh{generate.py}가 실패하더라도 마지막 \zsh{analyze.py}가 정상 종료하면 파이프라인 전체가 성공한 것처럼 보일 수 있다.

이 문제를 방지하기 위해 Bash와 zsh에서는 다음 옵션을 사용할 수 있다.

\begin{dhkzsh}
set -o pipefail
\end{dhkzsh}

이 옵션을 활성화하면 파이프라인에서 실패한 명령이 있을 때 그 실패를 파이프라인 종료 상태에 반영할 수 있다.

정확하게는 Bash와 zsh에서 \zsh{pipefail}을 사용하면 파이프라인의 종료 상태가 0이 아닌 종료 코드를 반환한 가장 오른쪽 명령의 종료 코드가 되고, 모든 명령이 성공하면 0이 된다.

ML/DL 자동화 스크립트에서 파이프라인을 사용할 때 유용한 설정이다.

\subsubsection{이번 학습 내용 정리}

이번에 배운 핵심 개념들을 정리하면 다음과 같다.

\begin{dhktable}
|개념|	의미|
|:---|:---|
|File Descriptor	|프로세스의 열린 입출력 자원을 식별하는 번호|
|0	|stdin|
|1	|stdout|
|2	|stderr|
|>	|stdout을 파일에 기록|
|>>	|stdout을 파일 끝에 추가|
|2>	|stderr를 파일에 기록|
|2>&1	|stderr의 출력 대상을 현재 stdout의 출력 대상과 동일하게 설정|
|<	|파일을 stdin에 연결|
|｜	|앞 명령의 stdout을 뒤 명령의 stdin에 연결|
|subprocess.PIPE	|Python에서 자식 프로세스의 입출력 파이프를 구성할 때 사용|
|pipefail	|파이프라인의 실패 상태를 더 정확하게 반영하는 Shell 옵션|
\end{dhktable}

지금 까지 배운 내용을 하나로 연결하면:

\begin{dhkdiagram}[in=4]
node os (1,2) :: Operating System
node sp (3,2) :: Shell Process
node ca (5,2) :: 명령 해석
node pa (6,1) :: Process A
node f0 (7,2) :: FD 0 (stdin)
node f1 (8,2) :: FD 1 (stdout)
node f2 (9,2) :: FD 2 (stderr)
node pi (8, 3) :: PIPE
node err (11,2) :: error.log
node pb (6,3) :: Process B
edge os -> sp
edge sp -- ca
edge ca -> pa :: HV
edge pa -- f0 :: VH
edge pa -- f1 :: VH
edge pa -- f2 :: VH
edge f2 -> err
edge f1 -- pi
edge pi -> pb
edge ca -> pb :: HV
\end{dhkdiagram}

이제 Shell이 단순히 명령을 입력받는 프로그램이 아니라, 프로세스를 실행하고 각 프로세스의 입출력 통로를 연결하는 역할도 수행한다는 것을 이해할 수 있을 것이다.

그리고 Python의 \zsh{print()}, \zsh{input()}, \zsh{sys.stdin}, \zsh{sys.stdout}, \zsh{sys.stderr}가 

왜 존재하는지도 더 명확해졌을 것이다.

다음 학습에서는 \zsh{subprocess.run()}과 \zsh{subprocess.Popen()}의 차이를 살펴볼 것이다.

\subsubsection{\zsh{subprocess.run()}과 \zsh{subprocess.Popen()}의 차이}

먼저 다음 코드를 보자.

\begin{dhkcode}
import subprocess
import sys

result = subprocess.run(
    [sys.executable, "-c", "print('Hello')"]
)

print("Finished")
\end{dhkcode}

이 코드는 자식 프로세스가 종료된 다음 \zsh{Finshed}를 출력한다.

실행 흐름은 다음과 같다:

\begin{dhksmry}[centered]
Parent Python Process
\da
subprocess.run()
\da
Child Process 생성
\da
Child Process 실행
\da
CHild Process 종료
\da
CompletedProcess 반환
\da
print("Finished")
\end{dhksmry}

그렇다면 자식 프로세스가 실행되는 동안 부모 프로세스에서 다른 작업을 하고 싶다면 어떻게 해야 할까?

이때 사용하는 것이 \zsh{subprocess.Popen()}이다.

\begin{dhkcode}
import subprocess
import sys

process = subprocess.Popen(
    [sys.executable, "-c", "print('Hello')"]
)

print("Parent is running")
\end{dhkcode}

\zsh{Popen()}은 자식 프로세스를 생성한 뒤,

그 프로세스가 종료될 때까지 기다리지 않고 \zsh{Popen} 객체를 반환한다.

따라서 부모 프로세스는 자식 프로세스가 실행 중인 상태에서도 다음 코드를 실행할 수 있다.

다만 두 프로세스의 출력 순서는 실행 스케줄링에 따라 달라질 수 있다.

두 함수의 차이

\begin{dhktable}
|구분	|subprocess.run()	|subprocess.Popen()|
|:---|:---|:---|
|주요 목적	|외부 프로그램 실행 후 결과 확인	|실행 중인 프로세스 직접 관리|
|기본 대기 동작	|종료까지 기다림	|생성 직후 반환|
|반환 객체	|CompletedProcess	|Popen|
|종료 코드	|result.returncode	|process.returncode|
|실행 중 상태 확인	|직접 지원하지 않음	|poll()|
|종료 대기|	내부적으로 수행|	wait()|
|입출력 통신	|input, capture_output 등	|communicate()|
|프로세스 종료 요청	|직접 제어에 제한적|	terminate(), kill()|
\end{dhktable}

여기서 중요한 사실이 하나 있다.

\zsh{subprocess.run()}은 내부적으로 \zsh{Popen()}을 이용해 구현된 고수준 함수다.

즉, 두 함수가 완전히 별개의 프로세스 실행 메커니즘을 사용하는 것은 아니다.

\zsh{run()}은 자주 사용하는 실행·대기·결과 수집 절차를 편리하게 묶어 놓은 인터페이스라고 이해하면 된다.

\subsubsection{\zsh{Popen} 객체와 프로세스의 관계}

이전 학습에서 Python의 이름과 객체 사이의 바인딩을 배웠다.

그 관점에서 다음 코드를 살펴보자.

\begin{dhkcode}
process = subprocess.Popen(
    [sys.executable, "-c", "print('Hello')"]
)
\end{dhkcode}

여기서 두 가지가 발생한다.

첫 번째는 운영체제 수준의 자식 프로세스 생성이다.

두 번째는 그 자식 프로세스를 관리하기 위한 Python의 \zsh{Popen} 객체 생성이다.

이를 구조로 표현하면 다음과 같다:

\begin{dhkbranch}[in=8]
node main :: Parent Python Process
branchline main -> name :: Python Namespace
line name
node name :: process
arrow name
node name :: Popen 객체
arrow name :: PID 등을 통해 관리
stop name
handoff main -> cp :: Child Process
line cp
node cp :: Python Interpreter
\end{dhkbranch}

여기서 반드시 구분해야 할 점이 있다.

\zsh{Popen} 객체 자체가 자식 프로세스인 것은 아니다.

\zsh{Popen} 객체는 부모 Python 프로세스의 메모리에 존재하는 Python 객체다.

반면 자식 프로세스는 운영체제가 별도로 관리하는 실행 단위이다.

따라서 다음 두 개념은 서로 다르다.

\begin{dhkzsh}
process
\end{dhkzsh}

이는 \zsh{Popen} 객체를 참조하는 Python 이름이다.

\begin{dhkzsh}
process.pid
\end{dhkzsh}

이는 자식 프로세스의 PID를 나타내는 속성이다.

예를 들어 다음 코드를 실행해 보자.

\begin{dhkcode}
import os
import subprocess
import sys

print("Parent PID:", os.getpid())

process = subprocess.Popen(
    [sys.executable, "-c", "import os; print(os.getpid())"]
)

print("Child PID:", process.pid)

process.wait()
\end{dhkcode}

예상 출력은 다음과 비슷하다.

\begin{dhkzsh}
Parent PID: 42001
Child PID: 42002
42002
\end{dhkzsh}

출력 순서는 달라질 수 있지만, 중요한 것은 부모와 자식의 PID가 서로 다르다는 사실이다.

또한 자식 프로세스가 종료되더라도 부모에 있는 \zsh{Popen} 객체가 자동으로 사라지는 것은 아니다.

Python 객체의 생명주기와 운영체제 프로세스의 생명주기는 별개의 개념이기 때문이다.

\subsubsection{\zsh{poll()} - 프로세스가 종료되었는지 확인하기}

\zsh{Popen()}은 자식 프로세스의 종료를 기다리지 않고 반환한다고 했다.

그렇다면 자식 프로세스가 아직 실행 중인지 어떻게 확인할까?

이때 사용하는 메서드가 \zsh{poll()}이댜.

\begin{dhkzsh}
process.poll()
\end{dhkzsh}

반환값은 다음과 같다.

\begin{dhktable}
|반환값|	의미|
|:---|:---|
|None|	아직 종료되지 않음|
|정수|	프로세스가 종료되었으며 해당 값은 종료 코드|
\end{dhktable}

예제를 보자.

\begin{dhkcode}
import subprocess
import sys
import time

process = subprocess.Popen(
    [
        sys.executable,
        "-c",
        "import time; time.sleep(3)"
    ]
)

print("Immediately:", process.poll())

time.sleep(4)

print("After 4 seconds:", process.poll())
\end{dhkcode}

예상 출력:

\begin{dhkzsh}
Immediately: None
After 4 seconds: 0
\end{dhkzsh}

여기서 \zsh{time.sleep(3)}은 자식 프로세스에서 3초 동안 실행을 지연시키는 코드다.

부모 프로세스에서는 4초를 기다렸으므로 일반적으로 자식 프로세스가 종료된 상태가 된다.

하지만 시스템 스케줄링 등에 따라 정확한 실행 시점은 달라질 수 있다.

\zsh{poll()}은 대기 함수일까?

아니다.

이 부분이 중요하다.

\zsh{poll()}은 현재 종료 상태를 확인하고 즉시 반환하는 메서드다.

자식 프로세스가 종료될 때까지 기다리지 않는다.

따라서 다음과 같은 코드를 작성할 수 있다.

\begin{dhkcode}
while process.poll() is None:
    print("Still running...")
    time.sleep(1)
\end{dhkcode}

이 코드는 자식 프로세스가 실행 중인 동안 주기적으로 상태를 확인한다.

이러한 방식을 \zsh{polling}이라고 부른다.

다만 너무 짧은 간격으로 반복 확인하면 불필요하게 CPU 자원을 사용할 수 있다.

\subsubsection{\zsh{wait()} - 프로세스 종료까지 기다리기}

이번에는 \zsh{wait()}을 살펴보자.

\begin{dhkzsh}
process.wait()
\end{dhkzsh}

\zsh{wait()}은 자식 프로세스가 종료될 때까지 기다린 뒤 종료 코드를 반환한다.

예제를 보자.

\begin{dhkcode}
import subprocess
import sys

process = subprocess.Popen(
    [
        sys.executable,
        "-c",
        "import time; time.sleep(3)"
    ]
)

print("Before wait")

returncode = process.wait()

print("After wait")
print("Return code:", returncode)
\end{dhkcode}

실행 결과:

\begin{dhkzsh}
Before wait
After wait
Return code: 0
\end{dhkzsh}

여기서 \zsh{After wait}은 자식 프로세스가 종료된 뒤에 출력된다.

\begin{dhknote}
\begin{dhksmry}[centered]
poll()과 wait()의 차이
\end{dhksmry}
\end{dhknote}

\begin{dhkbranch}
node main :: poll()
branchline main -> p :: 현재 상태 확인
stop p
handoff main -> r :: 즉시 반환
\end{dhkbranch}

\begin{dhkbranch}
node main :: wait()
branchline main -> p :: 자식 프로세스 종료 대기
stop p
handoff main -> r :: 종료 코드 반환
\end{dhkbranch}

이렇게 구분하면 된다.

그리고 \zsh{wait()}가 반환된 이후에는 다음 속성에서도 종료 코드를 확인할 수 있다.

\begin{dhkzsh}
process.returncode
\end{dhkzsh}

단 \zsh{Popen} 객체의 \zsh{returncode}는 프로세스가 종료되었다는 사실을

\zsh{poll()}, \zsh{wait()} 등의 방법으로 확인하기 전에는 \zsh{None}일 수 있다.

\subsubsection{\zsh{communicate()} - 자식 프로세스와 데이터 주고받기}

이제 중요한 메서드인 \zsh{communicate()}를 살펴보자.

우리는 앞에서 \zsh{subprocess.run()}을 사용할 때 다음과 같이 작성하였다.

\begin{dhkcode}
result = subprocess.run(
    command,
    capture_output=True,
    text=True
)
\end{dhkcode}

\zsh{Popen()}에서는 다음과 같이 작성할 수 있다.

\begin{dhkcode}
import subprocess
import sys

process = subprocess.Popen(
    [
        sys.executable,
        "-c",
        "print('Hello from child')"
    ],
    stdout=subprocess.PIPE,
    stderr=subprocess.PIPE,
    text=True
)

stdout, stderr = process.communicate()

print("stdout:", stdout)
print("stderr:", stderr)
print("returncode:", process.returncode)
\end{dhkcode}

예상 출력:

\begin{dhkzsh}
stdout: Hello from child

stderr:
returncode: 0
\end{dhkzsh}

여기에 새로운 코드가 등장했다.

\begin{dhkcode}
stdout=subprocess.PIPE
stderr=subprocess.PIPE
\end{dhkcode}

이는 자식 프로세스의 표준 출력과 표준 오류 출력을 파이프로 연결하겠다는 의미다.

그리고 다음 코드로 데이터를 수집한다.

\begin{dhkcode}
stdout, stderr = process.communicate()
\end{dhkcode}

\zsh{communicate()}는 자식 프로세스와 표준 입출력을 통해 통신하고, 프로세스 종료까지 기다린다.

반환값은 두 개의 요소로 구성된 튜플이다.

\begin{dhkzsh}
(stdout, stderr)
\end{dhkzsh}

따라서 튜플 언패킹을 통해 각각의 결과를 변수에 바인딩할 수 있다.

\begin{dhknote}
\begin{dhksmry}[centered]
communicate() 와 wait() 의 차이
\end{dhksmry}
\end{dhknote}

\zsh{wait()}는 기본적으로 프로세스 종료를 기다리는 역할이다.

반면 \zsh{communicate()}는 파이프를 통해 데이터를 주고받으면서 종료까지 기다리는 역할을 수행한다.

\begin{dhktable}
|메서드	|주요 역할|
|:---|:---|
|poll()	|종료 상태 확인|
|wait()|	종료 대기|
|communicate()	|표준 입출력 통신 및 종료 대기|
\end{dhktable}

여기서 한 가지 중요한 주의사항이 있다.

\zsh{stdout=PIPE} 또는 \zsh{stderr=PIPE}로 설정한 상태에서 자식 프로세스가 많은 데이터를 출력한다면,

단순히 \zsh{wait()}만 호출하는 것은 위험할 수 있다.

왜 그럴까?

\subsubsection{파이프 버퍼와 교착 상태(Deadlock)}

이 부분은 운영체제의 프로세스 모델을 이해하는 데 중요한 내용이다.

앞에서 파이프는 두 프로세스 사이에서 바이트 데이터를 전달하는 통로라고 하였다.

그런데 파이프에는 일반적으로 \zsh{유한한 크기의 버퍼(buffer)}가 존재한다.

다음 구조를 생각해 보자.

\begin{dhkbranch}
node main :: Child Process
arrow main :: stdout
node main :: PIPE BUFFER
arrow main
node main :: Parent Process
\end{dhkbranch}

자식 프로세스가 데이터를 출력하면 파이프 버퍼에 데이터가 쌓인다.

부모 프로세스가 그 데이터를 읽으면 버퍼에 공간이 생긴다.

그런데 부모가 데이터를 읽지 않는다면 어떻게 될까?

\begin{dhkbranch}[in=4]
node main :: Child Process
arrow main :: 계속 출력
node main :: PIPE BUFFER
arrow main :: 버퍼 가득 참
node main :: Parent Process
handoff main -> wait :: wait()로 자식 종료 대기
\end{dhkbranch}

버퍼가 가득 차면 자식 프로세스는 추가 데이터를 쓰기 위해 기다려야 할 수 있다.

반면 부모 프로세스는 자식 프로세스가 종료되기를 기다리고 있다.

결과적으로 다음과 같은 상태가 발생할 수 있다.

\begin{dhkdiagram}[in=4]
node pa (1,1) :: Parent
node ch (2,2) :: Child 종료를 기다림
node c (3,1) :: Child
node p (4,2) :: Pipe 버퍼 공간을 기다림
node r (6,2) :: Parent가 읽어야 함
edge pa -- ch :: VH
edge c -- p :: VH
edge r -> p
edge p -> ch
\end{dhkdiagram}

이것이 교착 상태의 한 예시이다.

따라서 파이프를 이용해 자식 프로세스의 출력을 수집할 때는

일반적으로 \zsh{communicate()}를 사용하는 것이 안전하다.

\zsh{communicate()}는 파이프에서 데이터를 읽으면서 자식 프로세스의 종료를 기다리도록 설계되어 있다.

다만 \zsh{communicate()}는 수집한 출력을 메모리에 보관하므로,

매우 큰 출력에는 적합하지 않을 수 있다.

이 점은 ML/DL 프로젝트의 로그 처리에서 다시 연결할 것이다.

\subsubsection{\zsh{Popen()}으로 두 프로세스를 실제로 연결하기}

이제 앞에서 배운 Shell 파이프를 Python으로 구현해 보자.

Shell에서는 다음과 같이 작성했다.

\begin{dhkzsh}
python generate.py | python analyze.py
\end{dhkzsh}

이를 \zsh{Popen()}으로 구현하면 다음과 같다.

\begin{dhkcode}
import subprocess
import sys

producer = subprocess.Popen(
    [sys.executable, "generate.py"],
    stdout=subprocess.PIPE
)

consumer = subprocess.Popen(
    [sys.executable, "analyze.py"],
    stdin=producer.stdout,
    stdout=subprocess.PIPE,
    text=True
)

producer.stdout.close()

output, _ = consumer.communicate()

producer_returncode = producer.wait()

print(output)
print("Producer:", producer_returncode)
print("Consumer:", consumer.returncode)
\end{dhkcode}

여기서 가장 중요한 부분은 다음이다.

\begin{dhkzsh}
stdout=subprocess.PIPE
\end{dhkzsh}

첫 번째 프로세스의 표준 출력을 파이프로 연결한다.

그리고 다음 코드에서

\begin{dhkzsh}
stdin=producer.stdout
\end{dhkzsh}

두 번째 프로세스의 표준 입력을 첫 번째 프로세스의 표준 출력 파이프에 연결한다.

전체 구조는 다음과 같다.

\begin{dhkbranch}[in=8]
node main :: Parent Python Process
branchline main -> prod :: producer
arrow prod
node prod :: generate.py
arrow prod :: stdout
node prod :: PIPE
arrow prod :: stdin
stop prod
handoff main -> con :: consumer
arrow con
node con :: analyze.py
arrow con :: stdout
node con :: PIPE
arrow con
node con :: Parent Python Process
\end{dhkbranch}

여기서 두 프로세스는 도잇에 실행될 수 있다.

첫 번째 프로세스가 데이터를 생성하면 두 번째 프로세스가 이를 읽어 처리할 수 있다.

\begin{dhknote}
\begin{dhksmry}[centered]
producer.stdout.close()는 왜 필요할까?
\end{dhksmry}
\end{dhknote}

다음 코드가 눈에 보인다.

\begin{dhkzsh}
producer.stdout.close()
\end{dhkzsh}

부모 프로세스가 가지고 있는 첫 번째 파이프의 읽기 쪽 핸들을 닫는 작업이다.

두 번째 프로세스가 이미 해당 파이프를 입력으로 전달받았으므로,

부모가 그 읽기 핸들을 계속 유지할 필요가 없다.

이렇게 하면 두 번째 프로세스가 조기에 종료되어 더 이상 파이프를 읽지 않을 때,

첫 번째 프로세스가 적절하게 파이프 단절을 감지하는 데 도움이 된다.

POOSIX 환경에서는 이 경우 생산자에게 \zsh{SIGPIPE}가 전달되거나 쓰기 작업이 실패할 수 있다.

중요한 것은 \zsh{close()}가 생산자 프로세스를 직접 종료시키는 명령은 아니라는 점이다.

파이프의 한쪽 끝에 대한 부모 프로세스의 참조를 닫는 것이다.

\subsubsection{\zsh{timeout} - 실행 시간 제한하기}

외부 프로그램을 실행할 때는 예상보다 오래 걸리는 상황이 발생할 수 있다.

예를 들어 모델 학습이 특정 단계에서 멈추거나, 외부 프로그램이 응답하지 않는 상황을 생각해 보자.

이때 사용할 수 있는 것이 \zsh{timeout}이다.

\begin{dhknote}
\begin{dhksmry}[centered]
subprocess.run()에서 사용하기
\end{dhksmry}
\end{dhknote}

\begin{dhkcode}
import subprocess
import sys

try:
    result = subprocess.run(
        [
            sys.executable,
            "-c",
            "import time; time.sleep(10)"
        ],
        timeout=3
    )

except subprocess.TimeoutExpired:
    print("Execution timed out")
\end{dhkcode}

이 코드는 자식 프로세스가 지정된 제한 시간 안에 종료되지 않으면

\zsh{TImeoutExpired} 예외를 발생시킨다.

\zsh{subprocess.run()}은 시간 초과 시 자식 프로세스를 종료시키고

종료 처리를 기다린 뒤 예외를 발생시키도록 구현되어 있다.

\begin{dhknote}
\begin{dhksmry}[centered]
\zsh{Popen()}에서 사용하기
\end{dhksmry}
\end{dhknote}

\zsh{Popen()}에서는 조금 다르다.

\begin{dhkcode}
import subprocess
import sys

process = subprocess.Popen(
    [
        sys.executable,
        "-c",
        "import time; time.sleep(10)"
    ]
)

try:
    process.wait(timeout=3)

except subprocess.TimeoutExpired:
    print("Timeout")

    process.terminate()
    process.wait()
\end{dhkcode}

여기서 중요한 차이가 있다.

\zsh{Popen.wait(timeout=...)}에서 시간 초과가 발생해도

자식 프로세스가 자동으로 종료되지는 않는다.

따라서 필요한 경우 직접 종료 요청을 해야 한다.

\begin{dhkzsh}
process.terminate()
\end{dhkzsh}

그리고 종료를 기다린다.

\begin{dhkzsh}
process.wait()
\end{dhkzsh}

이 차이는 실무에서 반드시 기억해야 한다.

\subsubsection{\zsh{terminate()}와 \zsh{kill()}}

프로세스가 실행 중일 때 종료를 요청하는 방법도 알아보자.

대표적으로 다음 두 메서드가 있다.

\begin{dhkzsh}
process.terminate()

process.kill()
\end{dhkzsh}

둘 다 프로세스 종료에 사용하지만 운영체제에 따라 동작 방식이 다르다

\begin{dhksmry}[centered]
macOS/Linux
\end{dhksmry}

POSIX 환경에서는 일반적으로 다음과 같이 동작한다.

\begin{dhktable}
|메서드	|동작|
|:---|:---|
|terminate()	|SIGTERM 신호 전달|
|kill()	|SIGKILL 신호 전달|
\end{dhktable}

\zsh{SIGTERM}은 프로세스에 종료를 요청하는 신호다.

프로세스가 이를 처리할 수 있기 때ㅔ문에 종료 전에 일부 정리 작업을 수행할 수도 있다.

반면 \zsh{SIGKILL}은 프로세스가 일반적인 방식으로 처리하거나 무시할 수 없는 강제 종료 신호다.

따라서 가능하다면 정상적인 종료 절차를 먼저 사용하는 것이 좋다.

\begin{dhksmry}[centered]
Windows
\end{dhksmry}

Windows에서는 내부적인 프로세스 종료 방식이 다르다.

특히 Python의 \zsh{Popen.terminate()}와 \zsh{Popen.kill()}은

Windows에서 모두 강제 종료에 가까운 동작을 수행한다.

따라서 macOS/Linux의 \zsh{SIGTERM}과 \zsh{SIGKILL} 구분을

Windows에 그대로 적용해서는 안 된다.

\begin{dhknote}
\begin{dhksmry}[centered]
terminate()를 호출하면 즉시 종료될까?
\end{dhksmry}
\end{dhknote}

반드시 그렇지는 않다.

\begin{dhkzsh}
process.terminate()
\end{dhkzsh}

는 종료를 요청하는 메서드다.

따라서 이후에

\begin{dhkzsh}
process.wait()
\end{dhkzsh}

를 호출해 실제 종료를 확인하는 것이 일반적이다.

또한 \zsh{terminate()}와 \zsh{kill()}은 기본적으로

해당 \zsh{Popen} 객체가 관리하는 프로세스를 대상으로 하며,

그 프로세스가 추가로 생성한 모든 자식 프로세스까지 자동으로 종료한다고 보장하지는 않는다.

\subsubsection{ML/DL 프로젝트에서 프로세스 제어하기}

이제 지금까지 배운 내용을 실제 프로젝트에 연결해 보자.

다음과 같은 구조가 있다고 가정하자.

\begin{dhkdir}
my-project/
> train.py
> run_experiment.py
> data/
>>trian.csv
>logs/
\end{dhkdir}

이번에는 학습 스크립트를 실행하면서 로그를 파일에 저장하고,

실행 상태를 확인하는 구조를 만들어 보자.

1. 학습 스크립트

\begin{dhkcode}
# train.py

import argparse
import time

parser = argparse.ArgumentParser()

parser.add_argument(
    "--epochs",
    type=int,
    default=5
)

args = parser.parse_args()

for epoch in range(1, args.epochs + 1):
    print(f"Epoch {epoch}/{args.epochs}", flush=True)

    time.sleep(1)

print("Training completed", flush=True)
\end{dhkcode}

여기서 \zsh{flush=True}가 새롭게 등장했다.

\begin{dhkcode}
print(..., flush=True)
\end{dhkcode}

일반적으로 표준 출력은 버퍼링될 수 있다.

특히 출력 대상이 터미널이 아니라 파일이나 파이프일 때는 출력이 즉시 전달되지 않을 수 있다.

\zsh{flush=True}는 \zsh{print()} 호출 후 해당 출력 스트림의 버퍼를 비우도록 요청하는 옵션이다.

따라서 학습 로그를 가능한 한 즉시 확인하고 싶을 때 유용하다.

2. 실험 실행 스크립트

\begin{dhkcode}
# run_experiment.py

import subprocess
import sys
import time
from pathlib import Path

project_root = Path(__file__).resolve().parent

train_script = project_root / "train.py"

log_dir = project_root / "logs"
log_dir.mkdir(parents=True, exist_ok=True)

log_path = log_dir / "train.log"

command = [
    sys.executable,
    str(train_script),
    "--epochs",
    "5"
]

with log_path.open("w", encoding="utf-8") as log_file:

    process = subprocess.Popen(
        command,
        cwd=project_root,
        stdout=log_file,
        stderr=subprocess.STDOUT
    )

    print("Training PID:", process.pid)

    while process.poll() is None:
        print("Training is running...")
        time.sleep(1)

    print("Training finished")
    print("Exit code:", process.returncode)
\end{dhkcode}

이 코드를 하나씩 분석해 보자.

\begin{dhknote}
\begin{dhksmry}[centered]
cwd=project_root
\end{dhksmry}
\end{dhknote}

\begin{dhkzsh}
cwd=project_root
\end{dhkzsh}

이 옵션은 \zsh{자식 프로세스의 현재 작업 디렉터리(CWD)}를 지정한다.

예전에 배운 \zsh{Path.cwd()}와 연결되는 부분이다.

부모 프로세스의 CWD를 변경하는 것은 아니다.

\begin{dhknote}
\begin{dhksmry}[centered]
stdout=log_file
\end{dhksmry}
\end{dhknote}

\begin{dhkzsh}
stdout=log_file
\end{dhkzsh}

자식 프로세스의 표준 출력을 파일에 연결한다.

Shell에서 다음 명령을 실행하는 것과 비슷하다.

\begin{dhkzsh}
python train.py > train.log
\end{dhkzsh}

\begin{dhknote}
\begin{dhksmry}[centered]
stderr=subprocess.STDOUT
\end{dhksmry}
\end{dhknote}

\begin{dhkzsh}
stderr=subprocess.STDOUT
\end{dhkzsh}

자식 프로세스의 표준 오류 출력을 표준 출력과 동일한 대상으로 연결한다.

Shell의 다음 표현과 유사한 의미이다.

\begin{dhkzsh}
python train.py > train.log 2>&1
\end{dhkzsh}

따라서 정상 출력과 오류 출력이 모두 \zsh{train.log}에 저장된다.

\begin{dhknote}
\begin{dhksmry}[centered]
process.poll()
\end{dhksmry}
\end{dhknote}

\begin{dhkcode}
while process.poll() is None:
\end{dhkcode}

자식 프로세스가 아직 실행 중인지 확인한다.

실행 중이면 부모 프로세스에서 상태 메시지를 출력한다.

로그 파일이 닫히는 시점

\begin{dhkcode}
with log_path.open("w", encoding="utf-8") as log_file:
\end{dhkcode}

\zsh{with} 블록이 끝나면 부모가 연 로그 파일 객체가 닫힌다.

여기서는 반복문이 자식 프로세스의 종료를 확인 한 뒤 \zsh{with} 블록을 빠져나가므로,

로그 파일을 너무 일찍 닫는 문제를 피할 수 있다.

\subsubsection{\zsh{subprocess.run()}과 \zsh{Popen()} 중 무엇을 선택할까?}

실무에서는 대부분의 단순한 외부 프로그램 실행에 \zsh{subprocess.run()}을 먼저 고려하면 된다.

예를 들어 다음처럼 실행 결과만 필요하다면:

\begin{dhkcode}
result = subprocess.run(
    command,
    capture_output=True,
    text=True,
    check=True
)
\end{dhkcode}

이 방식이 간결하고 관리하기 쉽다.

반면 다음과 같은 상황에서는 \zsh{Popen()}이 적합할 수 있다.

\begin{dhksmry}
- 외부 프로그램을 실행한 뒤 부모가 다른 작업을 계속해야 하는 경우

- 실행 중인 프로세스의 상태를 주기적으로 확인해야 하는 경우

- 여러 프로세스를 동시에 실행해야 하는 경우

- 프로세스의 표준 입출력을 직접 연결해야 하는 경우

- 실행 중인 프로세스에 종료 신호를 보내야 하는 경우
\end{dhksmry}

다만 \zsh{Popen()}이 더 저수준으 인터페이스이므로,

자식 프로세스의 종료 확인과 자원 정리를 개발자가 더 신경 써야 한다.

\subsubsection{이번 학습 내용 정리}

이번에 배운 핵심 개념은 다음과 같다.

\begin{dhktable}
|개념|	의미|
|:---|:---|
|Popen()	|자식 프로세스를 생성하고 관리 객체를 반환|
|Popen 객체|	자식 프로세스를 관리하는 Python 객체|
|pid	|자식 프로세스의 PID|
|poll()	|종료 상태를 확인하고 즉시 반환|
|wait()	|자식 프로세스 종료까지 대기|
|communicate()	|표준 입출력 통신과 종료 대기|
|subprocess.PIPE	|자식 프로세스의 표준 입출력을 파이프로 연결|
|subprocess.STDOUT	|stderr를 stdout과 동일한 대상으로 연결|
|timeout|	실행 시간 제한|
|terminate()	|자식 프로세스에 종료 요청|
|kill()|	자식 프로세스 강제 종료|
|cwd|	자식 프로세스의 작업 디렉터리 지정|
\end{dhktable}

전체 학습 구조는 다음과 같다

\begin{dhkbranch}[in=6]
node main :: Operating System
arrow main
node main :: Shell Process
arrow main :: python run_experiment.py
node main :: Parent Python Process
branchline main -> pi :: Python Interpreter
handoff pi -> exec :: sys.executable
stop exec
branchline main -> pe :: Process Environment 
branchline pe -> arg :: sys.argv 
stop arg
branchline pe -> env :: os.environ 
stop env
branchline pe -> cwd ::Path.cwd()
stop cwd
handoff pe -> getpid :: os.getpid()
stop getpid
handoff main -> sub :: subprocess.Popen()
branchline sub -> popen :: Popen 객체
branchline popen -> pid :: pid
stop pid
branchline popen -> poll :: poll()
stop poll
branchline popen -> wait :: wait()
stop wait
branchline popen -> comm :: communicate()
stop comm
branchline popen -> term :: terminate()
stop term
handoff popen -> kill :: kill()
stop kill
arrow sub
node sub :: Child process
branchline sub -> in :: stdin
stop in
branchline sub -> out :: stdout -> train.log
stop out
branchline sub -> err :: stderr -> train.log
stop err
handoff sub -> exit :: Exit code
arrow exit
node exit :: Popen.returncode
\end{dhkbranch}

여기서 가장 중요한 것은 Python 객체와 운영체제 프로세스를 구분하는 것이다.

\zsh{Popen} 객체는 부모 Python 프로세스에 존재하는 객체이고,

자식 프로세스는 운영체제가 별도로 관리하는 실행 단뒤이다.

그리고 \zsh{poll(), wait(), communicate()}는

그 자식 프로세스의 상태와 입출력을 관리하기 위한 인터페이스다.

이것으로 8장 프로세스와 외부 프로그램의 핵심 이론 학습은 완료했다.

물론 프로세스 간 통신(IPC), 시그널 처리, 프로세스 그룹, 멀티프로세싱 등 

더 깊이 공부할 수 있는 내용은 남아 있다.

하지만 현재 목표인 Python 프로젝트 실행 환경을 

이해하는 데 필요한 기본 범위는 충분히 다루었다.

이제 다음은 마지막 단계인 9장 — ML/DL 프로젝트에서 실제로 합치기다.

여기서는 지금까지 배운 개념을 개별적으로 설명하는 것이 아니라, 

하나의 프로젝트를 구성하면서 연결할 것이다.

예를 들어 다음 명령을 실행한다고 생각해 보자.

\begin{dhkzsh}
conda activate transformer

python -m src.train \
    --data ./data/train.csv \
    --epochs 10 \
    --batch-size 32 \
    --output ./checkpoints/model.pt
\end{dhkzsh}

이 명령을 중심으로 다음을 통합할 것이다.

\begin{dhksmry}
    - Conda 환경과 Python 인터프리터 선택

    - 프로젝트 루트와 CWD

    - python -m과 패키지 import

    - argparse와 Path

    - 환경변수와 .env

    - 데이터 파일 읽기

    - 모델 체크포인트 저장

    - 로그 디렉터리 관리

    - Hugging Face 모델 다운로드와 캐시

    - GPU/CPU 장치 선택

    - 학습 프로세스 실행과 종료 코드 관리
\end{dhksmry}
